\documentclass[a4paper,10pt]{article}
\input{glyphtounicode}
\pdfgentounicode=1
\usepackage[T1]{fontenc}
\usepackage{lmodern}
\usepackage[none]{hyphenat}
\usepackage[left=0.7in,right=0.7in,top=0.55in,bottom=0.55in]{geometry}
\usepackage[utf8]{inputenc}
\usepackage{enumitem}
\usepackage{titlesec}
\usepackage{hyperref}
\usepackage{xcolor}

\definecolor{darkblue}{RGB}{0,119,181}
\definecolor{linkedinblue}{RGB}{0,119,181}

\hypersetup{
    colorlinks=true,
    linkcolor=linkedinblue,
    urlcolor=linkedinblue,
    citecolor=linkedinblue,
    pdfborder={0 0 0}
}

\titleformat{\section}{\large\scshape\raggedright\color{darkblue}}{}{0em}{}[\titlerule]
\titleformat{\subsection}[runin]{\bfseries}{}{0em}{}[:]
\titlespacing*{\section}{0pt}{7pt}{4pt}
\titlespacing*{\subsection}{0pt}{5pt}{0.5em}

\setlength{\parindent}{0pt}
\sloppy
\setlist[itemize]{leftmargin=1.2em,itemsep=1pt,topsep=2pt,parsep=0pt}

\begin{document}
\pagestyle{empty}

\begin{center}
{\LARGE \textbf{Harvad Li}}\\[8pt]
{\Large \textit{AI Platform Engineer}}\\[8pt]
\textcolor{darkblue}{
  \href{mailto:hongzhili01@gmail.com}{hongzhili01@gmail.com} \textbar{}
  \href{tel:+46760067977}{+46 76 006 7977} \textbar{}
  Gothenburg, Sweden
}\\[3pt]
\textcolor{darkblue}{
  \href{https://www.linkedin.com/in/harvad/}{linkedin.com/in/harvad} \textbar{}
  \href{https://github.com/bluehawana}{github.com/bluehawana} \textbar{}
  \href{https://huggingface.co/bluehawana}{huggingface.co/bluehawana} \textbar{}
  \href{https://www.bluehawana.com}{www.bluehawana.com}
}
\end{center}

\section*{Profile}
Senior software engineer with 6+ years across full-stack development, cloud infrastructure and AI, now building and running the on-premises AI platform at Volvo Group Digital Technology \& Operations (DTO) and the AI agents that run on it. I own what I build end to end in production, including its roadmap, incidents and costs. I write code with AI coding agents (Claude Code) every day, and I know commercial-vehicle software from inside Volvo Group, one of Coretura's two owners.

\section*{AI at Volvo Group --- Highlights}
\begin{itemize}
  \item \textbf{2nd place of 16 teams at Volvo Group HackDay 2026}, leading the team behind ClaimGenerationBot; Volvo Trucks asked for a production rollout.
  \item \textbf{Built the first RAG system in Volvo Group DTO} (Support Chatbot, April 2026), in daily use by two teams and cutting support-ticket processing time by 35\%.
  \item \textbf{Four AI products in four months} (April to July 2026), each with a measured result: senior PR review time from 6 hours to 60 minutes, about 90\% less time to build a slide deck, and a three-agent system for trucks.
  \item \textbf{Run the platform underneath them:} open-weight models served on Volvo's own hardware, so no data goes to an outside AI service, at zero cloud API cost.
  \item \textbf{Teach it:} two company-wide sessions on RAG, agents and on-premises LLMs (74 and 71 registered participants).
\end{itemize}

\section*{Technical Skills}
\begin{itemize}
  \item \textbf{AI platform, model APIs and gateways:} Access to AI models across providers (Azure OpenAI, Anthropic Claude, Gemini API, Groq, NVIDIA NIM) and on-premises model serving (Ollama, MLX, SGLang, llama.cpp), AI gateway, OpenAI-compatible endpoints, failover and load balancing, model versions, benchmarking, quantization
  \item \textbf{AI agents and retrieval (RAG):} Agent harnesses and agent loops with tool calling, multi-agent systems, tool integrations, retrieval-augmented generation, knowledge bases, embeddings, vector search, prompts and prompt engineering, vision OCR, speech (Whisper, TTS)
  \item \textbf{Evaluations and observability:} Evaluations of answer quality (confidence scoring, grounded answers with sources, human escalation), model benchmarks, end-to-end test rigs, latency logging, audit trails, monitoring and alerting with Prometheus and Grafana
  \item \textbf{Security and safety:} Least privilege, secrets management (Azure Key Vault, managed identities, token rotation), Entra ID, data sovereignty, agent guardrails that limit prompt injection (tool allow-lists, grounded answers, human approval before actions)
  \item \textbf{Cloud infrastructure as code and CI/CD:} Terraform, Azure Bicep, Kubernetes (AKS, OpenShift ARO, on-premises), Helm, Argo CD, Docker, automated pipelines with Azure DevOps and GitHub Actions, resilience and failover, cost optimisation
  \item \textbf{Languages and testing:} Python (FastAPI, asyncio, pytest), C\#/.NET, TypeScript/React, Kotlin (Android Automotive, AOSP), SQL, PowerShell; well-tested, maintainable code, code review
  \item \textbf{AI coding agents:} Claude Code, GitHub Copilot, local coder models (Qwen Coder) in daily use; mentoring engineers to adopt them safely
\end{itemize}

\section*{Professional Experience}

\subsection*{Volvo Group --- Software Configurator, Software Factory, DTO, Gothenburg}
\textit{Full-Stack Developer \& AI Solution Engineer, Mar 2026 -- Present}
\begin{itemize}

  \item \textbf{Shared AI platform (on-premises):} Designed and run the model-serving layer behind the bots below: two Mac Studio M2 Ultra servers running Ollama with round-robin load balancing and automatic failover, exposed through OpenAI-compatible endpoints so each product can change model without code changes. Published 12+ custom models to \href{https://ollama.com/bluehawana}{Ollama} and \href{https://huggingface.co/bluehawana}{HuggingFace}, and contributed Apple Silicon concurrent-inference support to the open-source SGLang serving framework.

  \item \textbf{Volvo Smart Vehicle --- three agents on one truck record} (Jul 2026; Android Automotive/AOSP, Kotlin, FastAPI, Whisper, Gemma 4 on Ollama, MLX, Docker): Built a working prototype in which three agents read and write one record keyed by the registration plate and call each other instead of starting over. \textit{Co-Driver} is a hands-free voice agent in the cab with 33 tools that plans routes and rest breaks against EU driving-time rules, explains dashboard warning lamps from the handbook and books workshop service. \textit{ClaimBot} settles who pays for a repair from contract tier, warranty and open campaigns, and drafts the claim. \textit{Fleet Manager} gives the office a live view of every truck and relays customer changes to the driver. About 1 s per direct voice turn and 3--4 s when a turn crosses all three agents. Model chosen by benchmark against three alternatives, with an MLX failover. Agents only quote what the record returns and never report an unconfirmed booking; an end-to-end voice test rig and a 14-truck scenario fleet check every stage. Designed as one platform for any truck brand: the wake word and brand are configuration, and all data stays in the manufacturer's own cloud. Built with Claude Code. Code, architecture and a 10-minute demo: \href{https://github.com/bluehawana/CoDriver-LLM-TTS-Demo}{github.com/bluehawana/CoDriver-LLM-TTS-Demo}

  \item \textbf{ClaimGenerationBot --- 2nd place of 16 teams, Volvo Group HackDay 2026} (Python/FastAPI, .NET, RAG): Automated service-contract generation for Volvo Trucks: plate number in, the agent retrieves truck profile, service history and parts data, drafts the contract and books the workshop, and a human approves the final step. \textbf{From about 40 minutes to near-instant.}

  \item \textbf{Support Chatbot} (Apr 2026; .NET 10, SignalR, llama3.3:70b, nomic-embed-text, qwen3-vl, SQLite): Built the first RAG system in DTO, a support assistant on a Karpathy-style LLM wiki knowledge base for the Software Configurator and SWES teams, whose L1--L3 support engineers use it in their daily work. \textbf{Cut support-ticket processing time by 35\%.} Answers cite their sources and carry a confidence score, low-confidence answers are escalated to a specialist, screenshots are read with vision OCR, and every conversation is stored for audit.

  \item \textbf{DevBot --- AI pull-request review} (May 2026; Python, aiohttp, Microsoft Teams Bot Framework, Azure DevOps webhooks, Qwen Coder): Junior mode analyses a developer's changes and their risks and suggests how to get the PR approved. Senior mode shows what changed in the codebase and the likely result of approving the merge. \textbf{Cut average senior review time from 6 hours to 60 minutes.} Path-based review rules and incremental reviews, with documented secret handling and rotation.

  \item \textbf{OfficeBot --- slides from plain text} (Jun 2026; LLM, HTML): Turns plain text and a prompt into an HTML slide deck with the Volvo Group logo, colour theme and fonts applied automatically. \textbf{Saves about 90\% of the time} compared with building the deck by hand.

  \item \textbf{AI coding agents for the team, through the gateway} (PowerShell, Entra ID): Packaged a Claude Code setup that reaches Claude models through Volvo's GenAI Hub AI Gateway with short-lived Entra ID tokens instead of API keys, with an installer and a smoke test so any developer reproduces it in minutes.

  \item \textbf{AI assets as code, and security} (GitHub Actions, Azure Pipelines, Bicep, Azure Automation): Keep each bot's knowledge base and configuration in its own versioned repository, separate from the container image, so the application ships with no secrets and no Volvo data and a knowledge change is reviewed and traceable. CI and CD through GitHub Actions, with reusable Azure Pipelines templates for build, infrastructure, publish and security steps. Wrote a daily runbook that uses a managed identity with read-only Graph access to warn the team before app secrets and certificates expire.

  \item \textbf{Observability and resilience} (May 2026; Prometheus, Grafana, Alertmanager, HAProxy, .NET, Bicep): Replaced hard-to-read Splunk dashboards with Grafana dashboards the whole team reads at a glance: Argo CD sync status, pipeline runs, Azure DevOps activity per developer and a delivery health index built from recording rules. Added black-box probes that alert on the real cause when the service website is down or slow, and built a failover controller that moves traffic from on-premises hosts to Azure and back on Prometheus signals.

  \item \textbf{Team knowledge as data, and mentoring:} Wrote the team's coding handbook by reviewing every team repository, published as chapters for people and as a versioned knowledge base for retrieval, and built an LLM-generated 3D knowledge graph of the team's services and processes used for onboarding. Mentor colleagues adopting AI coding agents and local models.

\end{itemize}


\subsection*{ECARX Technology (Automotive Infotainment), Gothenburg}
\textit{Infrastructure Engineer \& Full-Stack Developer, Oct 2024 -- Nov 2025.} Cut cloud spend by 35\%, ran a zero-downtime migration from Azure AKS to on-premises Kubernetes (Helm, Argo CD, Terraform), and reduced MTTR by 35\% with Prometheus and Grafana.


\subsection*{Earlier roles}
.NET and Azure Integration Developer (Freelance), Synteda AB, Aug 2023 -- Sep 2024 \textbar{} Platform Architect \& Full-Stack Developer, Senior Material (Europe) AB, Jan 2022 -- Dec 2022



\section*{Open Source --- \href{https://github.com/bluehawana}{github.com/bluehawana}}
\textbf{\href{https://github.com/bluehawana/nvidia-nim-switch-python}{NVIDIA NIM Switch}}: model gateway that lets Claude Code run against 180+ NVIDIA NIM models, with a published compatibility report for all 182. \textbf{\href{https://github.com/bluehawana/Qwen3.827B-SGLang-mpbm5max}{Qwen 27B on Apple Silicon}}: OpenAI- and Anthropic-compatible gateway with admission control and metrics for concurrent serving.


\section*{Education, Certifications and Languages}
\textbf{IT-H\"ogskolan}, Gothenburg --- Bachelor's Degree in .NET Cloud Development, 2021--2023\\
\textbf{Campus M\"olndal} --- Java Integration, 2019--2021\\
\textbf{University of Gothenburg} --- Master's Degree in International Business and Trade, 2017--2019\\[3pt]
AWS Certified Solutions Architect -- Associate (2022) \textbar{} AWS Certified Data Analytics -- Specialty (2022) \textbar{} Microsoft Azure Fundamentals AZ-900 (2022) \textbar{} CNCF Scholarship Recipient (CKAD)
\\[3pt]
\textbf{Languages:} English (fluent) \textbar{} Swedish (B2, professional) \textbar{} Chinese (native)\\
\textbf{Work authorization:} Swedish permanent residence, based in Gothenburg


\end{document}
